\documentclass[../../main.tex]{subfiles}
\graphicspath{{../../image/}{../../figure/}}
\begin{document}

\section{注记}\label{section:ch18-注记}

关于二维空间中标量守恒律 Riemann 问题求解的早期与新近贡献，可见 Guckenheimer [2]、Wagner [1]、Lindquist [1]、Zhang 与 Zhang [1]、Tong Zhang 与 Yuxi Zheng [1]、Chen、Li 与 Tan [1] 以及 Xiaozhou Yang [1]。给出解的完整分类的详尽处理，包含在 Chang 与 Hsiao [3]、Li、Zhang 与 Yang [1] 以及 Yuxi Zheng [2] 的专著中。

气体流过各种形状的管道、喷管或障碍及其技术应用，早已被空气动力学家用解析、数值或实验方法研究。背景知识，包括此处记录在\cref{section:ch18-定常平面等熵气流}、\cref{section:ch18-自相似无旋等熵气流}中的定常或自相似平面无旋气流方程的设定，以及基本文献目录，可在任一本理论空气动力学标准教材中找到，例如 Courant 与 Friedrichs [1]、von Mises [1]、Hayes 与 Probstein [1] 以及 Ferrari 与 Tricomi [1]。一个非常有用而可靠的来源是 Serre [24] 的综述文章。另见 Serre [28] 对 von Neumann [5] 的评述。

1950 年代的严格分析最初集中于亚声速无旋流动，其中伪势所满足的方程为椭圆型；见 Shiffman [1] 与 Bers [1]。Morawetz [1,2,3] 开创了跨声速流动的研究，此时方程为椭圆--双曲混合型、在声速界面上退化。在经历三十年相对沉寂之后，由于椭圆或双曲型偏微分方程理论的进展，这一领域重新引起兴趣。Tai-Ping Liu [29] 给出了有趣的历史叙述。欲知这一正在进行的研究方向的现状，读者应参阅 Chen 与 Feldman [9] 即将出版的专著。综述与展望见 Gui-Qiang Chen [12]。

由于处理一般 Euler 方程组时遇到严重的解析困难，相当一部分研究是在无旋流动的较简单设定中、或在保留实际 Euler 方程组主要特征的简化方程组框架下进行的。例子有：压力梯度方程组，它由 Euler 方程组删去非线性对流项得到；Pressureless Euler 方程组，其中保留对流项但压强取为零；以及所谓小扰动方程组，通过启发式渐近论证得到。细节与详尽文献目录见 Yuxi Zheng [2] 的书与 Gui-Qiang Chen [11] 的综述文章。另一简化分析的途径是考察 Chaplygin 气体\eqref{align:ch02-2-5-23}，它使 Euler 方程组线性退化，从而相当大地简化了激波相互作用与波型；见 Serre [26,29]。

\cref{theorem:ch18-椭圆性原理}通常称为“椭圆性原理”，其证明见 Elling 与 Liu [1]。亚声速（椭圆）区域中关于压强的另一个有趣极大值原理，见 Serre [9]。

关于从双曲型向椭圆型过渡的另一视角，见 Serre 与 Freistühler [1]。

\cref{theorem:ch18-弱激波稳定性}取自 Chen、Zhang 与 Zhu [1]。另见 Chen、Xiao 与 Zhang [1]。\cref{theorem:ch18-Prandtl问题}给出了所谓 Prandtl 问题的解，归功于 Elling 与 Liu [2]，并由 Bae、Chen 与 Feldman [2] 改进。相关结果见 Gui-Qiang Chen 与 Beixiang Fang [1]、Gui-Qiang Chen 与 Tian-Hong Li [2]、Shuxing Chen [2,3,4,6]、Shuxing Chen 与 Beixiang Fang [1]、Shuxing Chen 与 Dening Li [1,2]、Chen、Min 与 Zhang [1]、Chen、Xin 与 Yin [1]、Lien 与 Liu [1]、Schaeffer [2]、Xin 与 Yin [2] 以及 Yongqian Zhang [1,2]。

关于\cref{section:ch18-激波撞凸壁}中概述的激波撞凸壁经典问题的确定理论解，见 Chen 与 Feldman [9]。特别地，这项工作证实了 von Neumann 的猜想：只要凸壁斜率超过脱体角，正规反射就是可能的。该方向的预备结果见 Chen 与 Feldman [4,6] 以及 Bae、Chen 与 Feldman [1]。同一设定下的进一步讨论见 Elling [1,2,3,4] 与 Sever [13]。曲壁上的激波碰撞见 Chen、Chen 与 Feldman [1]。在压力梯度方程组、小扰动方程组或 Chaplygin 气体的较简单设定下的正规反射，见 Čanić、Keyfitz 与 Kim [1,2,3]、Keyfitz 与 Lieberman [1]、Jegdić [1]、Jegdić、Keyfitz 与 Čanić [1]、Keyfitz [3]、Serre [26] 以及 Yuxi Zheng [3,4]。

有充分的数值与实验证据表明：当凸壁不够陡、不能支撑正规反射时，激波将脱开，形成 Mach 反射的波型；见 Ben-Dor [1] 的书。然而这一现象的严格理论处理仍付阙如。与 Mach 反射相关的问题，见 Čanić、Keyfitz 与 Kim [4]、Chen、Wang 与 Yang [1]、Shuxing Chen [8,9] 以及 Tesdall、Sanders 与 Keyfitz [1,2]。

喷管或管道中的跨声速流动，见于 Chen、Chen 与 Feldman [1]、Chen、Chen 与 Song [1]、Chen 与 Feldman [5,7,8]、Gui-Qiang Chen 与 Hairong Yuan [1]、Shuxing Chen [10]、Shuxing Chen 与 Hairong Yuan [1]、Xie 与 Xin [1,2,3]、Xin 与 Yin [1,3,4]、Bae 与 Feldman [1]、Chen、Deng 与 Xiang [1]、Tianyou Zhang 与 Yuxi Zheng [1]、Du、Weng 与 Xin [1]、Du、Xie 与 Xin [1]、Chunpeng Wang 与 Zhouping Xin [1]、Du、Xin 与 Yan [1] 以及 Li、Xin 与 Yin [1,2,3,4,5]。类似地，多维活塞问题由 Chen、Chen、Wang 与 Wang [1]、Shuxing Chen [5] 以及 Chen、Wang 与 Zhang [1,2] 讨论。

拐角处激波的衍射，见 Chen、Deng 与 Xiang [2]。

含稀疏波的全局无激波解，见 Heibig [3]、Kim 与 Song [1]、Jiequan Li [1,2]、Kyungwoo Song [1] 以及 Yuxi Zheng [5]。稀疏波的相互作用，研究于 Bang [1]、Glimm、Ji、Li、Zhang 与 Zheng [1]、Lei 与 Zheng [1]、Jiequan Li 与 Yuxi Zheng [1,2]、Li、Zhang 与 Zheng [1]、Li Yang 与 Zheng [1]、Chen 与 Zheng [1] 以及 Ji 与 Zheng [1]。

定常超声速二维气流的自相似小 $L^\infty$ 扰动结构，描述于 Elling 与 Roberts [1]。

二维空间中各类方程组 Riemann 问题的解，构造于 Čanić [1,2]、Čanić 与 Keyfitz [1,2]、Chang、Chen 与 Yang [1]、Schulz-Rinne [1]、Shuxing Chen [1]、Tan 与 Zhang [1]、Yang 与 Huang [1]、Zhang、Li 与 Zhang [1]、Tong Zhang 与 Yuxi Zheng [3] 以及 Yuxi Zheng [6]。

最后，关于跨声速激波与涡旋面的构造与稳定性，见 Chen 与 Feldman [1,2,3]、Shuxing Chen [7]、Chen、Zhang 与 Zhu [2]、Gui-Qiang Chen 与 Ya-Guang Wang [1]、Chen、Kukreja 与 Yuan [1,2]、Chen、Christoforou 与 Jegdić [1]、Keyfitz、Tesdall、Payne 与 Popivanov [1]、Mingjie Li 与 Yuxi Zheng [1] 以及 Colombel 与 Secchi [1,2,3,4]。

\cref{section:ch18-等距浸入}中讨论的、$\mathbb{R}^3$ 中曲面的等距浸入与无旋气体动力学之间奇特而有趣的联系，由 M. Slemrod 构想，并在 Chen、Slemrod 与 Wang [2,3,4] 中发展。另见 Christoforou [4] 与 Christoforou 与 Slemrod [1]。

非线性弹性中的空化最初由 Ball [2] 在静力学范围内处理。将该理论推广到弹性动力学（此处\cref{section:ch18-弹性动力学空化}中概述），由 Pericak-Spector 与 Spector [1,2] 开创，并由 Miroshnikov 与 Tzavaras [3] 进一步发展。关于如何压过空化失稳的有趣论证，见 Giesselmann 与 Tzavaras [1]。

关于气体动力学 Euler 方程的自相似径向对称解，已有大量文献；见 Courant 与 Friedrichs [1]、Yuxi Zheng [2]、Tong Zhang 与 Yuxi Zheng [4,5] 以及 Serre [10]。特别地，关于流体中显现空化的解，见 Wang 与 Li [1] 以及 Baisheng Yan [1]。

\medskip
\noindent\footnotesize 原书页 687 脚注：在该情形必须采用绝热气（非等熵）气体动力学\eqref{align:ch03-Euler}的完整方程组，因为等熵方程组\eqref{align:ch03-等熵Euler}退化为平凡。

\end{document}
